\documentclass[a4paper,11pt]{article}
\usepackage[T1]{fontenc}
\usepackage{lmodern}
\usepackage{amsmath,amssymb,amsthm}
\usepackage[margin=1in]{geometry}
\usepackage{url}
\usepackage{microtype}
\usepackage[hidelinks]{hyperref}
\hypersetup{pdftitle={An Explicit GCD Evaluation for Three-Direction Torus Arrangements},pdfauthor={Rushikesh Murlidhar Walke},pdfsubject={AI-assisted expository note; author-review copy v08; CC BY 4.0; not peer-reviewed},pdfkeywords={toric arrangements, cyclic grids, Euler characteristic, arithmetic Tutte polynomial}}
\newtheorem{theorem}{Theorem}
\newtheorem{lemma}{Lemma}
\newtheorem{corollary}{Corollary}
\newcommand{\Torus}{\mathbb{R}^2/\mathbb{Z}^2}
\title{An Explicit GCD Evaluation for Three-Direction Torus Arrangements}
\author{Rushikesh Murlidhar Walke\\
\small Independent Researcher, Maharashtra, India\\
\small ORCID: \url{https://orcid.org/0009-0006-6749-5660}}
\date{Author-review copy v08, 6 November 2026}

\begin{document}
\maketitle

\begin{center}
\small AI-assisted expository note. Not peer-reviewed.
\par\smallskip
\footnotesize \copyright\ 2026 Rushikesh Murlidhar Walke.
This article is licensed under
\href{https://creativecommons.org/licenses/by/4.0/}{CC BY 4.0}.
\par Original supplementary software: MIT License; see \path{anc/LICENSE}.
\end{center}

\begin{abstract}
We explicitly evaluate established toric-arrangement counting principles
for a family of closed curves on the two-dimensional torus. The vertical
coordinates form the union of two common-origin cyclic grids of orders
$N$ and $R$; the other two families are the level sets of $x+y$ and $x-y$
on grids of orders $Y$ and $K$. A finite-group calculation gives the number
of triple intersections, and a cellularity argument justifies the use of
Euler's relation. For $R=12$, $Y=27$, and $K=60$, the face count is
$87N+4248-84\gcd(N,12)-3\gcd(N,1080)$. In particular, it is $6264$ at
$N=27$. The result concerns geometric faces and does not require a sampled
trajectory or numerical experiment. We give an elementary proof and
identify the same expression as an arithmetic-Tutte specialization.
\end{abstract}

\section{Relation to existing work}
Chamber counting for toric arrangements is established theory.
Ehrenborg, Readdy and Slone give a general count from the intersection
poset under an open-ball hypothesis~\cite[Theorem 3.6]{ERS}.
For two-dimensional arrangements, the identity
$F=\sum_{j\geq2}(j-1)t_j$, with $t_j$ counting intersections of $j$ curves,
appears in the preprint version of Chandrasekhar and
Deshpande~\cite[Lemma 2.4]{CD}. Our $F=P-T$ is its double/triple-point case.
Moci's arithmetic multiplicities and compact-torus chamber
formula~\cite[Section 2.2, Lemma 5.4 and Corollary 5.16]{Moci}
also yield the expression evaluated here.

The purpose of this note is to display the gcd reduction for the specified
common-origin grids, including duplicate vertical levels, and provide a
self-contained proof and parameterized example. Neither Euler's relation
nor a new general chamber-count principle is claimed as a contribution.
The precise literature connections are made explicit below; finding no
identical printed numerical example would not establish priority.

The author's earlier expository note~\cite[pp. 2--5, Theorem 1 and
Corollary 1]{Walke2026} counts pairs on finite cyclic grids subject to
sum and difference constraints, using Smith normal form. Its formula
specializes exactly to the triple-point count used here, as shown in
Section~\ref{sec:earlier}. The present note adds the torus arrangement,
the cellularity argument, and the passage from point incidences to
geometric faces. The two notes therefore share a finite-grid counting
component while counting different final objects. Neither note claims a
new general counting theory.

\section{Setup and scope}
For a positive integer $m$, let
\[
 G_m=\{j/m+\mathbb Z:0\leq j<m\}\subset\mathbb R/\mathbb Z.
\]
All grids have origin zero; independently rotated grids are outside the
stated assumptions. Fix positive integers $N,R,Y,K$. On $\Torus$ consider
the curves
\[
 x=u\quad(u\in U:=G_N\cup G_R),\qquad
 x+y=a\quad(a\in G_Y),\qquad
 x-y=b\quad(b\in G_K),
\]
where equations are understood modulo one. A vertical curve occurring in
both grids is included only once. Each individual curve is an embedded
circle. Curves in the same family are disjoint, and curves from different
families meet transversely.

Vertices are distinct intersection points. Edges are the arcs between
successive vertices on each curve, allowing a loop edge when there is only
one vertex on a curve. Faces are connected components of the complement.
Write their counts as $V,E,F$, and write $T$ for the number of points on
three curves. There are no points on more than three distinct curves.

The problem is one of combinatorial topology. Euler's relation below is
the surface formula for a graph embedding, not the Euler equations of
fluid mechanics. We make no claim here about orbital sampling, error
detection, astronomical accuracy, or the number of distinct labels after
a relabelling of the faces.

\section{Face-count theorem}
Put
\[
 g=\gcd(N,R),\quad v=N+R-g,\quad d=\gcd(Y,K),\quad
 L=\operatorname{lcm}(Y,K),\quad P=v(Y+K)+2YK.
\]

\begin{theorem}\label{thm:main}
For the arrangement just defined, every face is an open disk. Moreover,
\begin{align}
 T&=d\bigl(\gcd(N,2L)+\gcd(R,2L)-\gcd(g,2L)\bigr),\label{eq:triple}\\
 V&=P-2T,\qquad E=2P-3T,\qquad F=P-T.\label{eq:counts}
\end{align}
These formulas hold for all positive integers $N,R,Y,K$.
\end{theorem}

\begin{lemma}\label{lem:triple}
If the vertical family is a single grid $G_n$, its triple-point count is
\[
 TT(n;Y,K)=d\gcd(n,2L).
\]
\end{lemma}
\begin{proof}
The homomorphism
\[
 G_Y\times G_K\longrightarrow G_L,\qquad (a,b)\longmapsto a+b
\]
is surjective. Indeed, $L/Y=K/d$ and $L/K=Y/d$ are coprime, so Bezout's
identity shows that $1/L$ belongs to $G_Y+G_K$. Thus this sum is $G_L$.
Its kernel has size $YK/L=d$, so each element of $G_L$ has exactly
$d$ preimages.

A fixed vertical coordinate $u$ supports a triple point exactly when
$a+b=2u$ for some $(a,b)\in G_Y\times G_K$. Each such pair gives the
unique point $(u,a-u)$, and distinct pairs give distinct points at that
vertical coordinate. Thus there are $d$ triple points at $u$ if
$2u\in G_L$, and none otherwise. The condition $2u\in G_L$ is equivalent
to $u\in G_{2L}$. Finally,
$G_n\cap G_{2L}=G_{\gcd(n,2L)}$, which proves the claim.
\end{proof}

\begin{lemma}\label{lem:cellular}
The faces of the arrangement are open disks.
\end{lemma}
\begin{proof}
Lift all the curves to affine lines in $\mathbb R^2$. Each chamber of the
lifted complement is an intersection of open strips, hence is convex.
In particular it lies between two consecutive vertical lines and between
two consecutive lines of the form $x+y=\text{constant}$. These two strip
directions are independent, so the chamber is bounded. The line
arrangement is locally finite, so each chamber is a bounded convex open
polygon and is homeomorphic to an open disk.

If a chamber intersects one of its integer translates, the two are the
same connected component of the lifted complement. A nonzero translation
cannot preserve a bounded nonempty chamber. Consequently the covering
projection identifies no two points of its interior. Its image is a
whole component of the torus complement: a path in that component lifts
to a path remaining in the chamber. Thus each face is homeomorphic to a
lifted chamber. Boundary identifications can occur and do not invalidate
cellularity of this graph embedding.
\end{proof}

\begin{proof}[Proof of Theorem~\ref{thm:main}]
Since $G_N\cap G_R=G_g$, there are $v$ distinct vertical curves. A vertical
curve meets each sum curve once and each difference curve once. A pair
$x+y=a$, $x-y=b$ has exactly two solutions: the equation $2x=a+b$ has
two roots modulo one, and each root determines $y$. Therefore the total
pairwise intersection count, with multiplicity, is $P$.

Apply Lemma~\ref{lem:triple} to $G_N$, $G_R$, and their intersection $G_g$.
Inclusion--exclusion gives equation~\eqref{eq:triple}. If $D$ denotes
the number of double points, then $P=D+3T$ and $V=D+T$, giving
$V=P-2T$.

Double and triple vertices have degrees four and six respectively. Every
curve has a vertex because the other two families are nonempty. The
degree sum, counting a loop twice, gives
$2E=4D+6T$, hence $E=2P-3T$. Lemma~\ref{lem:cellular} makes this a finite
cellular embedding on the torus. Euler's relation $V-E+F=0$ now gives
$F=E-V=P-T$.
\end{proof}

\section{Connection with the earlier cyclic-grid note}\label{sec:earlier}
For comparison with~\cite[Eq. (1)]{Walke2026}, write
\[
 C(a,b;Y,K)=\#\{(u,h)\in G_a\times G_b:
                  u+h\in G_Y,\ u-h\in G_K\}.
\]
The earlier note's union count is
\[
 Q(a,b;Y,K;r)=C(a,b;Y,K)+C(r,b;Y,K)
                         -C(\gcd(a,r),b;Y,K).
\]
Here $d=\gcd(Y,K)$ and $L=\operatorname{lcm}(Y,K)$ retain the
notation of this note; the earlier note uses different letters for some
of these quantities.

At a triple point, $2y=(x+y)-(x-y)\in G_L$, so $y\in G_{2L}$.
Consequently, restricting the second coordinate to $G_{2L}$ loses no
triple points, and there are exact identifications
\[
 TT(n;Y,K)=C(n,2L;Y,K),\qquad T=Q(N,2L;Y,K;R).
\]
Conversely, each pair in the specified $C$ or $Q$ set is a triple point
of the corresponding arrangement, so these identifications are bijective.

To display the specialization of the earlier Smith-normal-form formula,
put $q=\gcd(n,2L)$, $n=q a$, $2L=q b$, and
$M=\operatorname{lcm}(n,2L)=qab$. Then $\gcd(a,b)=1$, and the
integer matrix in~\cite[Eqs. (2)--(4)]{Walke2026} is
\[
 A=\begin{pmatrix}Yb&Ya\\Kb&-Ka\end{pmatrix}.
\]
Its invariant factors are
\[
 s_1=d,\qquad s_2=\frac{2YK\,ab}{d}=2L\,ab=Mb.
\]
Since $d\mid M$, the earlier formula gives
\[
 C(n,2L;Y,K)
 =\frac{n(2L)}{M^2}\gcd(M,s_1)\gcd(M,s_2)
 =\frac{n(2L)}{M}d=dq.
\]
Its inclusion--exclusion formula~\cite[Corollary 1, Eq. (8)]{Walke2026}
therefore gives exactly equation~\eqref{eq:triple}. This is an algebraic
specialization of the earlier exposition, not a separate originality
claim for the triple-count arithmetic. The earlier note does not count
complementary faces of a torus arrangement; the geometric argument above
is needed to pass from $T$ to $F$.

\section{An arithmetic-Tutte specialization}
This section is a direct specialization of the cited general results,
not a further general counting theorem. Use the indexed list
\[
 X=((N,0),(R,0),(Y,Y),(K,-K)).
\]
Its character kernels give exactly the union of curves studied here.
For a rank-two sublist $A$, its arithmetic multiplicity $m(A)$ is the
gcd of its nonzero absolute two-by-two determinants. The region count
is the arithmetic-Tutte evaluation
\[
 M_X(1,0)=\sum_{\substack{A\subseteq X\\r(A)=2}}
             (-1)^{|A|-2}m(A)\qquad\text{\cite{Moci}}.
\]
Here $r(A)$ is the rank over $\mathbb Q$ of the indexed character sublist.
Subsets are subsets of indexed positions, even if $N=R$.
Write $S(n)=\gcd(nY,nK,2YK)=d\gcd(n,2L)$. The following evaluation is
our algebraic substitution into that known formula:
\begin{align*}
 \text{pairs:}&\quad (N+R)(Y+K)+2YK,\\
 \text{triples:}&\quad -g(Y+K)-S(N)-S(R),\\
 \text{four positions:}&\quad +S(g).
\end{align*}
Their sum is $P-T$. Thus the explicit formula is already a consequence
of established arithmetic-arrangement theory. The elementary proof above
avoids needing that framework to verify this particular evaluation.

\section{The parameters \texorpdfstring{$R=12$, $Y=27$, $K=60$}{R=12, Y=27, K=60}}
Here $d=3$ and $L=540$, so
\[
 TT(n;27,60)=3\gcd(n,1080).
\]
Equivalently, writing $\nu_p(n)$ for the exponent of a prime $p$ in $n$,
\[
 TT(n;27,60)=2^{\min(\nu_2(n),3)}
             3^{\min(\nu_3(n)+1,4)}
             5^{\min(\nu_5(n),1)}.
\]
This equivalence follows from $1080=2^3\cdot3^3\cdot5$; it does not
depend on extrapolating a finite numerical pattern. The function $TT$
itself is not multiplicative in the standard arithmetic sense, since
$TT(1)=3$. Its normalization $TT(n)/3=\gcd(n,1080)$ is multiplicative.

\begin{corollary}
For every positive integer $N$, with the other three parameters fixed
as above,
\[
 \boxed{F(N)=87N+4248-84\gcd(N,12)-3\gcd(N,1080).}
\]
\end{corollary}
\begin{proof}
Set $g=\gcd(N,12)$. Since $g$ divides both $12$ and $1080$, the triple
count is $T=3\gcd(N,1080)+36-3g$. Also
$P=87(N+12-g)+3240$. Subtracting proves the result.
\end{proof}

At $N=27$ one obtains
\[
 v=36,\quad T=108,\quad P=6372,\quad
 (V,E,F)=(6156,12420,6264).
\]
For the immediate integer neighbours,
\[
 F(26)=6336,\qquad F(27)=6264,\qquad F(28)=6336.
\]
Thus $N=27$ is a strict local minimum relative to these neighbours.
It is not a local maximum. These statements do not assert any physical
optimality or lack of arithmetic structure.

For clarity, $TT(28;27,60)=12$ counts triples on the single grid $G_{28}$,
whereas $T=36$ counts triples on $G_{28}\cup G_{12}$. These quantities
must not be interchanged.

\section{What the count does not identify automatically}
The number of geometric faces need not equal the number of distinct
tuples produced by labelling them. For example, $N=R=Y=K=1$ gives
$T=1$, $P=4$, and $(V,E,F)=(2,5,3)$, while the four one-bin raw labels
are constant. A claim equating faces with a particular labelled encoding
requires a separate injectivity or connectedness argument. A further
claim that a finite trajectory visits every face requires additional
analysis. Neither identification is used in Theorem~\ref{thm:main}.

\section{Supplementary implementation}
The companion file \texttt{anc/hv27\_torus\_counts.py} provides the gcd
formula and a rational vertex-enumeration method using only Python's
standard library. The latter counts incidences at exact rational vertices
and applies the cellular Euler relation justified above. Comparing the
two implementations is a consistency check; the all-parameter theorem
is established by the proof, not by such comparisons.

The accompanying \texttt{anc/run\_hv27\_checks.py} defines a finite
validation suite. Author-supplied reports dated 6 November 2026 record a
Windows CPython 3.14.7 run with 284 completed checks, all passing: two
reference examples, 26 default-parameter cases with $15\leq N\leq40$,
eight nondefault-parameter examples, 200 pure-grid triple-count comparisons
with $1\leq n\leq200$, and 48 invalid-input cases. The pure-grid test uses
the integer congruence
\[
 K(2iY-jn)\equiv0\pmod{nY},\qquad 0\leq i<n,\quad 0\leq j<Y,
\]
instead of the companion's rational vertex construction.

The supplied reports are included as \path{anc/validation/RESULTS.json}
and \path{anc/validation/RESULTS.txt}. Their recorded checks agree, and
the recorded core and runner SHA-256 hashes match the distributed source
files. These are inspected author-supplied runtime records, not an
independent re-execution by the preparation assistant. They do not record
a terminal process exit code. Earlier archived logs concern other source
versions and are not used to validate this release.

These finite checks support implementation consistency. In particular,
the vertex enumerator obtains $F$ from $E-V$; it does not independently
enumerate the complement's faces. The cellularity proof and the
all-parameter theorem remain mathematical arguments, rather than
conclusions drawn from a finite run.

\clearpage
\section*{Assistance and author review}
ChatGPT (OpenAI) and Claude (Anthropic) were used for AI assistance in
the research and preparation of this note. Across the workflow, this
assistance included calculations, proof arguments, mathematical exposition,
literature comparison, and supplementary source code. It was not limited to language
editing. Historical computation records and the new author-supplied
validation reports were read as data. No research code or proof checker
was executed by the preparation assistant.

The author formulated the motivating mathematical idea and initial
hypotheses, supplied the starting basis for the investigation, and pursued
the research with AI assistance. The reported results were developed
through this internally AI-assisted process. The author also supplied the
publication identity and new runtime reports. This contribution statement
records the author's account of his role in originating and developing
the project; it makes no priority claim for the established mathematical
methods used here.

Internal mathematical investigation and inspection were AI-assisted.
Completed personal verification of every proof step is not asserted, and
no independent human peer review of this draft is documented. AI tools
are not authors.

\begin{thebibliography}{9}
\raggedright
\bibitem{CD}
K. Chandrasekhar and P. Deshpande,
\emph{Face enumeration for line arrangements in a 2-torus},
Indian Journal of Pure and Applied Mathematics \textbf{48} (2017), 345--362.
\url{https://doi.org/10.1007/s13226-017-0234-7}.
Preprint consulted: \url{https://arxiv.org/abs/1404.1665v1}.
Lemma numbers cited here refer to the preprint.

\bibitem{ERS}
R. Ehrenborg, M. Readdy and M. Slone,
\emph{Affine and toric hyperplane arrangements},
Discrete \& Computational Geometry \textbf{41} (2009), 481--512.
\url{https://doi.org/10.1007/s00454-009-9134-x}.
Preprint consulted: \url{https://arxiv.org/abs/0810.0295v1}.

\bibitem{Moci}
L. Moci,
\emph{A Tutte polynomial for toric arrangements},
Transactions of the American Mathematical Society \textbf{364}(2)
(2012), 1067--1088.
\url{https://doi.org/10.1090/S0002-9947-2011-05491-7}.
Preprint consulted: \url{https://arxiv.org/abs/0911.4823v5}.
Result numbers cited here refer to the preprint version consulted.

\bibitem{Walke2026}
R. M. Walke,
\emph{Counting intersections of rational cyclic grids: an application
of Smith normal form}, version 9.1, Zenodo preprint, 6 October 2026.
\url{https://doi.org/10.5281/zenodo.23184983}.
Full seven-page text consulted in the author-supplied PDF;
Theorem 1 and Corollary 1 are used in the comparison above.
\end{thebibliography}

\end{document}
